% Chapter 6: Análisis del Cliente y Segmentación
\chapter{Análisis del Cliente y Segmentación}

\section{Segmentación de Clientes Operadores}
Los clientes de la consultoría de ingeniería minera se dividen en tres segmentos principales con comportamientos de compra radicalmente diferentes:

\begin{enumerate}
    \item \textbf{Majors (Grandes Operadores Globales):} BHP, Anglo American, Codelco, Vale, Freeport-McMoRan.
    \begin{itemize}
        \item \textbf{Comportamiento:} Procesos de licitación altamente formales y prolongados (6-12 meses). Priorizan la reputación global, la capacidad financiera (balance sheet) de la consultora y el cumplimiento estricto de estándares globales.
        \item \textbf{Necesidades:} Grandes contratos EPCM de infraestructura, consultoría estratégica de descarbonización y soporte en megaproyectos.
        \item \textbf{Sensibilidad al Precio:} Media-Baja; valoran más la mitigación de riesgos catastróficos que el costo de la ingeniería.
    \end{itemize}

    \item \textbf{Mid-Tier (Productores Medianos):} Lundin Mining, Capstone Copper, Antofagasta Minerals (en ciertos activos), Polyus, Seligdar.
    \begin{itemize}
        \item \textbf{Comportamiento:} Decisiones más rápidas y pragmáticas. Valoran la agilidad, la experiencia específica en la geografía local y la capacidad de proponer diseños eficientes en CAPEX.
        \item \textbf{Necesidades:} Estudios de optimización de plantas, ampliación de depósitos de relaves, optimización de flota minera.
        \item \textbf{Sensibilidad al Precio:} Media-Alta; buscan la mejor relación costo-valor técnico.
    \end{itemize}

    \item \textbf{Juniors (Exploradoras y Desarrolladores de Proyectos):}
    \begin{itemize}
        \item \textbf{Comportamiento:} Altamente dependientes de financiamiento externo (mercado de capitales). Decisiones muy rápidas.
        \item \textbf{Necesidades:} Estudios de Factibilidad Definitivos (DFS) para levantar capital.
        \item \textbf{Sensibilidad al Precio:} Extrema; el costo del estudio consume su caja directa de exploración.
    \end{itemize}
\end{enumerate}

\section{Propuesta de Valor de REDCO por Segmento}
Para una firma como REDCO, el foco óptimo de captura son los operadores \textbf{Mid-Tier} y el soporte especializado a \textbf{Majors} a nivel regional (particularmente en Chile y Perú). La propuesta de valor de REDCO se estructura de la siguiente manera:

\begin{marketdatabox}[Matriz de Alineación de Valor REDCO]
\begin{itemize}[leftmargin=*]
    \item \textbf{Para Majors (Nivel Regional):} Proporcionar soporte técnico flexible y ágil para destrabar cuellos de botella en operaciones existentes (Ingeniería de Contraparte / Optimización), evitando la burocracia de los gigantes de EPCM.
    \item \textbf{Para Mid-Tiers:} Ingeniería de optimización de CAPEX (estilo Ausenco) pero con profunda sensibilidad regulatoria y comunitaria local (Chile/Perú).
    \item \textbf{Para Juniors:} Estudios NI 43-101 eficientes y rápidos para acelerar el levantamiento de capital en las bolsas de Toronto y Londres.
\end{itemize}
\end{marketdatabox}
